%% 第三章

\chapter{深空场景三维重建与新视角合成问题分析}
\chaptermark{深空场景三维重建与新视角合成问题分析}

\section{深空场景数据特征与方法适用性观察}

在分析现有三维重建与新视角合成方法在深空场景中的退化原因之前，有必要首先从数据层面观察深空场景与地面场景之间的条件差异。这里的差异并不只是图像外观上的远近、明暗或颜色不同，而是体现在深度分布、纹理结构、有效光度监督以及跨视角光照一致性等多个方面。这些数据特征会直接影响多视角几何约束的可靠性、图像重建损失的监督分布以及基于光度一致性的优化过程。

图3-1给出了地面场景与深空场景在图像内容和相对深度分布上的对比。地面场景样例中，主要物体和背景之间虽然存在一定深度变化，但整体深度分布相对集中，图像中的可辨识结构和局部纹理也较为明确。相比之下，深空场景样例中同时包含近处岩石、坡面结构和远处大范围背景区域，辅助深度结果显示出更明显的远近跨度。

\begin{figure}[H]
    \centering
    \begin{subfigure}[b]{0.48\linewidth}
        
\includegraphics[width=\linewidth]{fig3_1a_regular_scene_panel.png}
        \caption{地面场景图像}
    \end{subfigure}\hfill
    \begin{subfigure}[b]{0.48\linewidth}
        
\includegraphics[width=\linewidth]{fig3_1b_regular_depth_panel.png}
        \caption{地面场景相对深度}
    \end{subfigure}
    \par\medskip
    \begin{subfigure}[b]{0.48\linewidth}
        \includegraphics[width=\linewidth]{fig3_1c_deepspace_scene_panel.png}
        \caption{深空场景图像}
    \end{subfigure}\hfill
    \begin{subfigure}[b]{0.48\linewidth}
        \includegraphics[width=\linewidth]{fig3_1d_deepspace_depth_panel.png}
        \caption{深空场景相对深度}
    \end{subfigure}
    \caption{地面场景与深空场景的图像及相对深度分布对比。地面场景图像来源于Unsplash，深空场景图像来源于NASA/JPL-Caltech/ASU/MSSS。}\label{fig:3-1}
\end{figure}

由图3-1可以看出，深空场景中的远近结构差异更加突出。这意味着同一幅图像中的近处目标、坡面和远处背景并不处在相同的观测条件下，它们在像素尺度、纹理清晰度和视差信息上存在差异。远区域因此不能简单视作背景，而是会影响三维表示完整性和新视角合成质量的关键部分。

除深度跨度外，深空场景图像还常表现出低纹理、低对比和局部结构重复等特点。以火星沙丘或大面积砂面区域为例，图像中大量砂地和远处背景呈现缓慢变化的亮度分布，明显边缘和强纹理主要集中在少数沙丘边界、岩石轮廓或明暗过渡位置。虽然这些低纹理区域同样参与全图光度损失计算，但其局部颜色变化较弱，能够提供的有效梯度信息有限。图3-2从原始图像、梯度响应和高梯度结构区域叠加三个角度展示了这种光度监督分布不均衡现象。

\begin{figure}[H]
    \centering
    \begin{subfigure}[t]{0.32\linewidth}
        \centering
        \includegraphics[width=\linewidth]{fig3_2a_original.png}
        \caption{原始图像}
    \end{subfigure}\hfill
    \begin{subfigure}[t]{0.32\linewidth}
        \centering
        \includegraphics[width=\linewidth]{fig3_2b_gradient_response.png}
        \caption{梯度响应}
    \end{subfigure}\hfill
    \begin{subfigure}[t]{0.32\linewidth}
        \centering
        \includegraphics[width=\linewidth]{fig3_2c_high_gradient_overlay.png}
        \caption{高梯度结构区域叠加}
    \end{subfigure}
    \caption{深空场景图像中的光度监督不充分现象。}\label{fig:3-2}
\end{figure}

由图3-2可以看出，深空沙丘场景中存在大面积砂地、弱纹理、重复石块以及远处低对比区域。梯度响应图显示，较强的图像梯度主要集中在少数石块轮廓、沙丘棱线、沙纹边缘和阴影边界附近；高梯度结构区域叠加图进一步说明，强光度约束集中于有限的边界和明暗过渡处，未标出的大片砂面区域则属于相对弱约束区域。这说明全图光度损失虽然覆盖整幅图像，但有效约束主要来自少数结构位置，大面积砂面和远处低对比区域提供的监督更弱。

跨视角光照变化是深空场景中的另一类重要数据特征。图3-3展示了同一深空场景在不同太阳直射角条件下的图像及亮度分布差异。两幅图中的地形结构具有较强对应关系，但阴影边界、局部高亮区域和整体灰度分布发生了明显变化。

\begin{figure}[H]
    \centering
    \begin{subfigure}[b]{0.48\linewidth}
        \includegraphics[width=\linewidth]{fig3_3a_light_angle1_panel.png}
        \caption{光照条件一}
    \end{subfigure}\hfill
    \begin{subfigure}[b]{0.48\linewidth}
        \includegraphics[width=\linewidth]{fig3_3b_light_angle2_panel.png}
        \caption{光照条件二}
    \end{subfigure}
    \par\medskip
    \begin{subfigure}[b]{0.82\linewidth}
        \includegraphics[width=\linewidth]{fig3_3c_brightness_hist.png}
        \caption{亮度直方图}
    \end{subfigure}
    \caption{同一深空场景在不同太阳直射角条件下的光度差异。上方两幅图展示太阳直射角变化导致的明暗、阴影和局部高亮差异，下方亮度直方图展示对应图像的亮度分布变化。}\label{fig:3-3}
\end{figure}

由图3-3可以看出，不同观测图像之间的差异并不完全来自几何结构或纹理细节，而包含由成像光照条件变化引起的光度分布偏移。亮度直方图中的分布差异说明，同一场景在不同太阳直射角下会呈现不同的明暗比例和高亮区域范围，这会削弱多视角光度一致性假设，使同一空间区域在训练中接收到不完全一致的颜色监督。

综合图3-1至图3-3可以看出，深空场景的困难来自深度、纹理和光照条件的共同作用。深度跨度改变远近区域的可观测性，低纹理和重复结构使强光度监督集中于少数边界，跨视角光照变化则削弱颜色一致性。

基于上述观察，本文将现有方法在深空场景中面临的问题归纳为三类：其一，极端深度跨度使远处结构和背景层次在全图损失中贡献不足；其二，弱纹理、低对比和有限视角使光度优化难以充分约束高斯尺度和空间覆盖范围，容易诱发基元形态退化；其三，跨视角光照不一致会引入颜色监督偏差，影响外观表达稳定性。后续各节将围绕这三类问题展开分析。


\section{极端深度跨度导致的远区域监督不足}

上一节已经说明，深空场景中的近处结构与远处背景往往具有显著深度差异。本节进一步从3DGS优化过程分析这种深度跨度带来的影响：当远区域在图像中呈现更小尺度、更弱纹理和更低梯度时，它虽然参与损失计算，却更难产生足够有效的参数更新。

基础3DGS的训练通常依赖渲染图像与真实图像之间的全图光度重建误差。该误差对整幅图像统一计算，但优化器实际接收的是由像素误差反向传播得到的梯度信号。近处边缘、岩石轮廓和明显纹理更容易产生较大误差变化；远处地表的细碎结构在像素尺度上更小，颜色变化也更缓，容易被全图平均损失稀释。因此，全图损失覆盖所有像素，并不等价于所有深度区域获得同等强度的有效监督。

图3-4进一步以同一测试视角中的远处ROI为例，展示深度位置与局部渲染细节之间的对应关系。红框区域位于图像左上区域向右下偏移的远处地表，包含小石块边缘、浅暗纹理变化和细碎地形起伏；下排比较GT和不同方法在该ROI内的放大结果。

\begin{figure}[H]
    \centering
    \begin{subfigure}[t]{0.48\linewidth}
        \centering
        \includegraphics[width=\linewidth]{fig3_4a_gt_roi_box.png}
        \caption{GT测试视角与ROI}
    \end{subfigure}\hfill
    \begin{subfigure}[t]{0.48\linewidth}
        \centering
        \includegraphics[width=\linewidth]{fig3_4b_aux_depth.png}
        \caption{辅助深度图与ROI}
    \end{subfigure}
    \par\medskip
    \begin{subfigure}[t]{0.24\linewidth}
        \centering
        \includegraphics[width=\linewidth]{fig3_4d_gt_roi.png}
        \caption{GT ROI}
    \end{subfigure}\hfill
    \begin{subfigure}[t]{0.24\linewidth}
        \centering
        \includegraphics[width=\linewidth]{fig3_4e_3dgs_roi.png}
        \caption{3DGS ROI}
    \end{subfigure}\hfill
    \begin{subfigure}[t]{0.24\linewidth}
        \centering
        \includegraphics[width=\linewidth]{fig3_4f_3dhgs_roi.png}
        \caption{3D-HGS ROI}
    \end{subfigure}\hfill
    \begin{subfigure}[t]{0.24\linewidth}
        \centering
        \includegraphics[width=\linewidth]{fig3_4g_mvgs_roi.png}
        \caption{MVGS ROI}
    \end{subfigure}
    \caption{远区域监督不足导致的局部细节退化可视化。第一行展示测试视角图像及其辅助深度图，红框标出了用于局部放大的相对远处ROI。辅助深度图中颜色越亮表示相对深度越大。}\label{fig:3-4}
\end{figure}

由图3-4下排可以看出，GT图像中的远处ROI仍然包含可观察的小尺度结构，例如碎石边缘、细碎纹理以及局部明暗起伏。相比之下，3DGS和3D-HGS在该区域的渲染结果中出现明显模糊化，部分小石块边界变软，颗粒状纹理和浅暗变化被削弱，远处地表的细节对比度下降。MVGS虽然相对保留了更多局部结构，但其远处细碎纹理和边缘清晰度仍弱于GT。该结果说明，远区域并非缺少真实结构，而是这些细小结构在全图光度优化中容易被弱化，导致新视角渲染结果出现细节保留不足的问题。

从方法假设角度看，问题的根源在于基础3DGS默认由全图光度误差自动平衡不同区域的重要性，但深空数据中远区域的初始化、匹配和光度梯度往往同时偏弱。若训练目标不区分深度区域，优化会更倾向于降低近处和强纹理区域误差，远处局部结构则可能被平滑或低频化。

因此，远区域退化本质上是深度跨度与全图平均监督机制之间的不匹配。后续方法需要显式识别并强化远区域约束，使模型对远处边缘、碎石纹理和背景层次形成更充分的表达。

\section{光度优化约束不足导致的高斯基元形态退化}

3DGS将场景表示为一组三维高斯基元，每个高斯通常包含位置、尺度、旋转、不透明度和颜色特征等参数。其中，尺度参数直接决定高斯在三维空间中的形态和覆盖范围，是影响场景表示稳定性的重要因素。在观测视角充分、图像纹理丰富、几何初始化可靠的情况下，多视角光度约束能够在一定程度上限制高斯基元的空间形态，使其逐步形成较合理的分布。

但在深空场景中，上述条件往往难以同时满足。一方面，真实探测数据的视角数量和基线范围受到任务轨迹、观测窗口和传感器姿态限制，难以提供均匀充分的环绕视角。另一方面，地表图像中常存在大面积弱纹理、重复纹理或阴影区域，使特征匹配和多视图几何约束变弱。此时，仅依赖投影后的图像重建误差来优化三维高斯参数，可能无法充分约束每个高斯在真实三维空间中的尺度形态。

在上述光度约束不足的情况下，高斯基元形态退化通常表现为三类情况。第一，某一尺度方向显著大于其他方向，使高斯退化为细长或针状结构；第二，多个尺度方向同时增大，使高斯覆盖过宽空间区域，导致局部边界被平滑；第三，局部高斯尺度分布突变，造成相邻区域之间空间支撑关系不连续。这些异常形态未必在正常训练视角渲染中立即暴露，因为较大的投影覆盖和颜色累积可能暂时掩盖三维支撑的不合理性，但在新视角或诊断渲染中更容易表现为纹理拖拽、结构断裂和局部模糊。

为衡量高斯基元的各向异性程度，本文采用三个尺度方向中的最大值与最小值之比作为尺度比例：
\[
r_s=\frac{s_{\max}}{s_{\min}}
\]
其中，\(s_{\max}\) 和 \(s_{\min}\) 分别表示单个高斯基元三个尺度方向中的最大尺度和最小尺度。尺度比例越大，说明该高斯在某一方向上的拉伸越明显，其空间形态越偏离均衡分布。

图3-5从图像渲染、覆盖诊断和参数统计三个层面展示了基础3DGS中的形态退化现象。GT图像用于说明真实场景应有的地表结构和纹理状态；基础3DGS正常渲染结果用于观察常规渲染下的局部外观退化；Scaling Modifier = 0.5 的覆盖诊断可视化通过缩小高斯投影覆盖范围，使原本被较大覆盖区域掩盖的空间支撑不连续更容易显现；尺度比例分布则从参数统计层面反映训练后高斯基元的各向异性程度。

\begin{figure}[H]
    \centering
    \begin{subfigure}[t]{0.48\linewidth}
        \centering
        \includegraphics[width=\linewidth]{fig3_5a_gt.png}
        \caption{GT图像}
    \end{subfigure}\hfill
    \begin{subfigure}[t]{0.48\linewidth}
        \centering
        \includegraphics[width=\linewidth]{fig3_5b_3dgs_render.png}
        \caption{基础3DGS正常渲染结果}
    \end{subfigure}
    \par\medskip
    \begin{subfigure}[t]{0.48\linewidth}
        \centering
        \includegraphics[width=\linewidth]{fig3_5c_scale_modifier.png}
        \caption{Scaling Modifier = 0.5的覆盖诊断可视化}
    \end{subfigure}\hfill
    \begin{subfigure}[t]{0.48\linewidth}
        \centering
        \includegraphics[width=\linewidth]{fig3_5d_scale_ratio_distribution.png}
        \caption{高斯尺度比例分布}
    \end{subfigure}
    \caption{光度优化约束不足导致的高斯基元形态退化可视化。GT图像展示真实地表结构，基础3DGS渲染结果展示常规渲染下的局部退化现象；缩小高斯投影覆盖后，覆盖缺失和空间支撑不连续更容易暴露；尺度比例分布进一步说明部分高斯基元存在明显各向异性。}\label{fig:3-5}
\end{figure}

由图3-5可以看出，基础3DGS在正常渲染下虽然能够恢复主要地表外观，但局部弱纹理区域仍容易出现细节平滑和结构连续性不足。进一步将Scaling Modifier设为0.5后，部分区域出现明显的覆盖空洞、线状拖拽和支撑不连续现象，说明正常渲染中较大的高斯覆盖可能掩盖了内部空间形态的不稳定。尺度比例分布也呈现出明显长尾，表明训练后存在一部分尺度比例较大的高斯基元，这些基元具有较强各向异性，容易在新视角下引入不稳定投影。

从优化原理上看，基础3DGS的图像重建损失主要约束最终渲染图像与真实图像之间的像素差异，而不是直接约束三维高斯的尺度比例是否合理。只要某些异常高斯在训练视角下能够通过投影、遮挡和颜色混合降低图像误差，它们就可能被保留下来。对于纹理充分、视角覆盖较好的数据，多视角一致性能够在一定程度上抑制这种退化；但在深空场景中，弱纹理、低对比、局部结构重复和视角覆盖不足会削弱这种隐式限制。

因此，深空场景中的高斯基元形态退化问题，本质上来自基础3DGS光度优化目标对三维尺度形态的显式约束不足。仅依赖二维图像重建误差很难保证高斯尺度分布稳定，需要在不破坏3DGS显式可微渲染框架的前提下，引入针对尺度比例和空间覆盖范围的辅助约束。

\section{跨视角光照条件不一致导致的光度优化偏差}

多视角三维重建与新视角合成方法通常隐含一个基本假设：同一三维区域在不同视角下具有相对一致的光度外观。对于传统多视图几何而言，该假设支撑特征匹配和光度一致性；对于3DGS而言，该假设也影响颜色参数、透明度和几何位置的联合优化。当同一空间结构在不同图像中呈现相近颜色和亮度时，模型更容易将不同视角的观测融合为稳定的三维表示。

深空场景中的跨视角外观条件却往往并不一致。不同图像可能受到太阳高度角、阴影范围、传感器曝光、白平衡偏移、局部反射差异以及压缩噪声等因素影响。同一岩石、坑壁、斜坡或目标表面在不同视角下可能出现亮度增强、局部变暗、颜色偏移或阴影边界变化。这些变化并不一定代表真实几何或材料发生改变，却会直接影响基于光度误差的训练过程。

当模型缺少外观变化建模能力时，它可能将成像条件变化错误解释为场景本身属性变化。例如，某一视角中的阴影可能被模型吸收到高斯颜色中，导致其他视角下渲染偏暗；过曝区域可能使模型降低局部纹理对比度，造成细节丢失；白平衡偏移可能使不同视角的颜色监督互相冲突，最终产生平均化的颜色结果。这类问题会降低新视角渲染的外观一致性，并可能进一步干扰几何和透明度优化。

在现有方法直接应用于此类深空场景数据时，上述光度差异会表现为更具体的渲染局限。基础3DGS或其他依赖光度一致性假设的方法，可能在强阴影边界、局部高亮区域和弱纹理区域出现颜色漂移、暗部细节缺失或局部过平滑等问题。这些现象并不是单一模型容量不足造成的，而是与跨视角观测中不一致的颜色监督直接相关。

\begin{quote}
\textbf{配图提示：}此处应放现有方法面对深空场景数据时的局限可视化图，建议包含真实图像、基础3DGS及若干对比方法在同一局部区域的渲染结果，用于展示阴影、高亮或弱纹理区域中的颜色偏移、细节损失和局部伪影。
\end{quote}

从优化角度看，跨视角光照不一致会使同一高斯在不同视角中接收到相互冲突的颜色监督。基础3DGS通常使用固定或视角相关的球谐颜色特征表示外观，但这种表示主要用于建模观察方向带来的颜色变化，并不必然能够区分场景固有外观和成像条件变化。当图像间存在较强曝光或颜色偏移时，模型可能通过折中方式拟合多视角观测，从而导致所有视角下都不够准确。

若优化目标仍将这些差异统一视为由同一组颜色参数解释的重建误差，模型就会在多个观测之间进行折中，最终可能表现为暗部偏色、高亮区域细节损失或局部颜色漂移。因此，跨视角光照不一致会从监督信号层面改变颜色优化方向，而不仅仅是在渲染结果中产生表面性的亮度差别。

由此可见，光照条件不一致并不是单纯的图像预处理问题，而是会直接破坏多视角光度一致性假设，并影响3DGS的颜色、透明度和局部结构优化。对于深空场景，如果希望模型在复杂外观条件下保持稳定的新视角合成能力，就需要在颜色预测或外观表达层面补充对跨视角外观变化的建模。

\section{深空场景核心挑战归纳}

综合上述分析，常规三维重建与新视角合成方法在深空场景中的退化并不是由单一因素造成的，而是由数据特点与方法优化机制之间的不匹配共同导致。极端深度跨度会使远区域在全图损失中被削弱；低纹理、低对比和局部结构重复会削弱光度误差对高斯尺度形态的约束，而基础3DGS缺少显式形态正则会进一步放大高斯基元退化问题；跨视角光照条件不一致则会使光度监督本身产生偏差。这三类问题分别作用于图像监督、三维表示和颜色外观三个层面，因而需要从不同角度进行针对性设计。

\begin{table}[H]
    \fontsize{10.5pt}{13.125pt}\selectfont
    \caption{深空场景中的主要问题及对应方法需求}
    \centering\label{tab:3-1}
    \begin{adjustbox}{max width=\textwidth}
    \begin{tabular}{cccc}
        \toprule
        问题现象 & 原因分析 & 方法需求 & 对应设计 \\
        \midrule
        远区域细节弱化 & 全图损失被近处或强纹理区域主导 & 增强远区域监督 & FAME \\
        高斯基元形态退化 & 光度优化缺少显式尺度形态约束 & 稳定尺度分布与空间覆盖 & RVS \\
        颜色偏移与外观不一致 & 跨视角光度一致性假设受破坏 & 建模视角间外观变化 & 外观一致性建模 \\
        \bottomrule
    \end{tabular}
    \end{adjustbox}
\end{table}

表3-1展示了第三章问题分析与后续方法设计之间的对应关系。可以看到，深空场景中的主要挑战并不是简单地要求模型具有更强表达能力，而是要求模型能够在监督分布、显式几何表示和外观建模三个层面适应数据特点。对于远区域监督不足，需要使远处像素在训练中获得更明确的优化信号；对于高斯基元形态退化，需要从尺度方向分布和整体体积变化角度约束高斯基元；对于跨视角光照不一致，需要在颜色预测过程中区分静态外观和动态成像条件。

因此，本文后续提出的DSS-GS方法并不是对基础3DGS的简单模块堆叠，而是针对深空场景中暴露出的三类关键问题进行设计。第四章将围绕FAME远区域增强、RVS高斯尺度形态稳定以及跨视角外观一致性建模三个部分展开，说明本文方法如何从图像监督、三维表示和外观表达三个层面改善深空场景的新视角合成质量。

\section{本章小结}

本章围绕常规三维重建与新视角合成方法在深空场景中的适用性问题进行了分析。首先通过定性现象说明基础3DGS及部分主流方法在深空场景中可能出现远区域模糊、局部结构不稳定和颜色偏移等问题；随后从全图损失监督不均衡、光度优化对高斯尺度形态约束不足和跨视角光度一致性失效三个角度解释问题产生原因；最后归纳出远区域监督增强、高斯尺度形态稳定和外观一致性建模三项方法需求，为第四章DSS-GS方法设计提供依据。
